\documentclass[11pt]{article}

\usepackage{acl}

\usepackage{times}
\usepackage{latexsym}
\usepackage[T1]{fontenc}
\usepackage[utf8]{inputenc}
\usepackage{microtype}
\usepackage{inconsolata}
\usepackage{graphicx}

\title{Replication and Draw-Aware Extension Plan for Football Outcome Prediction with Classical Machine Learning}

\author{
Benjamin Tran 
\AND 
Nicholas Ramirez-Ornelas \\[0.5em]
California State Polytechnic University, Pomona \\
CS 5170 Advanced Natural Language Processing
}

\begin{document}

\maketitle

\begin{abstract}
This project replicates \textit{Forecasting Football Results Using Machine Learning Techniques} \cite{armanda2025forecasting}, which predicts football match outcomes as a three-class classification task: home win, draw, or away win. The original work evaluates Random Forest, XGBoost, and AdaBoost under both imbalance and SMOTE-balancing. We implement the prediction pipeline with preprocessing, feature engineering, class balancing, model training, and evaluation using accuracy, macro-F1, ROC-AUC, and confusion matrices. Our replication finds that model performance is sensitive to dataset composition, preprocessing choices, and balancing strategy. While ensemble models are able to achieve moderate accuracy, draw prediction is still the most difficult class because draws are less frequent and often occur when teams are closely matched in skill. Compared with the original paper, our results align with the broader finding that tree-based models are effective for football prediction, but differences in exact rankings highlight the importance of reproducibility controls. As an extension, we propose a draw-aware approach using parity-based uncertainty features, cost-sensitive training, and probability calibration to improve draw-class F1 without substantially reducing the overall accuracy.
\end{abstract}

\section{Introduction}

Predicting football match outcomes remains a difficult multi-class classification problem because of outcome imbalance and the uncertainty associated with factors such as team form, player performance, and tactical decisions. In particular, draw outcomes are consistently more difficult to predict than home or away wins due to their lower frequency and higher variability. Recent machine learning approaches have applied ensemble and boosting models to improve prediction accuracy in football forecasting tasks.

The replicated source paper, \textit{Forecasting Football Results Using Machine Learning Techniques} \cite{armanda2025forecasting}, studies this problem using Random Forest, XGBoost, and AdaBoost models under both imbalance and SMOTE-balancing. The paper reports moderate prediction accuracy improvements while also identifying draw prediction as the most challenging outcome class across models. 

Our project contribution is threefold: (1) a clean and reproducible reimplementation pipeline covering preprocessing, feature engineering, balancing, training, and evaluation; (2) transparent replication-gap reporting that analyzes differences between our findings and the original paper; and (3) a draw-aware extension focused on improving draw-class prediction performance through parity-based uncertainty features, cost-sensitive learning, and probability calibration.


\section{Related Work}

Football outcome prediction has been studied using both statistical and machine learning approaches. Recent work has shown that machine learning models, particularly ensemble and boosting methods, can improve prediction accuracy for sports forecasting tasks. Bunker and Thabtah \cite{bunker2019sport} proposed a machine learning framework for sports result prediction, while Decroos et al. \cite{decroos2019actions} highlighted the importance of contextual football event analysis. 

Ensemble learning methods such as Random Forest and XGBoost are used because of their ability to model nonlinear relationships in tabular sports data. Random Forest was introduced by Breiman \cite{breiman2001random}, while Chen and Guestrin \cite{chen2016xgboost} developed XGBoost as an efficient gradient boosting framework which is optimized for scalability and predictive accuracy.

Class imbalance still remains a challenge in football prediction, especially for draw outcomes. Chawla et al. introduced SMOTE for minority oversampling \cite{chawla2002smote}, while Sjoberg \cite{sjoberg2023match} discussed machine learning approaches for football match prediction. Wright \cite{wright2020} explored the integration of qualitative and contextual information in sports analytics to improve predictive models. Van Buuren \cite{vanbuuren2018flexible} further discussed methods for handling missing data during preprocessing. Our work extends this literature through a replication-focused analysis and a draw-aware extension focused on improving draw-class prediction performance.

\section{Background}

The original paper develops the football outcome prediction as a three-class classification problem consisting of home win, draw, and away win outcomes. The study uses football match data collected from Kaggle and additional match related sources, creating a dataset containing over 12,500 matches and more than 200 features. These features include team statistics, match information, and contextual variables such as home or away conditions.

The preprocessing pipeline includes missing-value handling, outlier removal, one-hot encoding, min-max normalization, and feature selection using Recursive Feature Elimination (RFE) with a Random Forest classifier. Numerical values are imputed using the median, categorical values use mode imputation, and outliers are capped at the 95th percentile. The paper selects the top 20 most important features for prediction.

For model development, the paper evaluates Random Forest, XGBoost, and AdaBoost. The Random Forest model uses 100 decision trees with a maximum depth of 10, while XGBoost uses a learning rate of 0.1 and 150 estimators with early stopping to reduce overfitting. The dataset is split into 80\% training and 20\% testing data, and a 5-fold cross-validation is used for evaluation.

The models are evaluated using accuracy, precision, recall, F1-score, ROC-AUC, and confusion matrices. The paper also compares both imbalance and SMOTE-balanced training settings to address the class imbalance, particularly for draw predictions. The results indicated that Random Forest with SMOTE balancing achieved the best overall performance.

\section{Replication Methodology}

\subsection{Pipeline}

Our replication reimplements the preprocessing, feature engineering, model training, and evaluation pipeline described in the original paper. The pipeline is organized as a configuration-driven experiment framework using a YAML configuration file (configs/replication.yaml) to ensure reproducibility across the runs. The implementation supports both imbalanced and SMOTE-balanced training settings for a direct comparison with the original study.

The preprocessing pipeline follows the methodology described in the source paper. Numerical features use median imputation for missing values, while categorical features use mode imputation. Outliers are capped using winsorization at the 95th percentile. Categorical variables are transformed through one-hot encoding, and numerical features are normalized using min-max scaling. Feature selection is performed using a Random Forest based selector to retain the top 20 predictive features.

Three machine learning models are evaluated: Random Forest, XGBoost, and AdaBoost. The Random Forest model uses 100 estimators with a maximum depth of 10, while XGBoost uses 150 estimators with a learning rate of 0.1. AdaBoost is configured with 200 estimators and a learning rate of 0.8. The implementation supports optional SMOTE oversampling to address class imbalance, particularly for draw outcomes.


\subsection{Reproducibility Controls}

To improve reproducibility, the replication framework uses configuration controlled experiment execution with fixed random seeds (seed = 42) across preprocessing, feature selection, model training, and cross-validation. The dataset is split using either temporal or stratified training test partitioning, depending o n the configuration settings.

The system automatically exports the evaluation metrics, confusion matrices, selected feature lists, and configuration snapshots for each experiment run. In addition, 5-fold cross-validation is performed to evaluate the model consistency across the data subsets. This reproducibility-focused design allows controlled comparisons between balanced and imbalanced training settings while reducing the experimental variability.

\section{Replication Results}

Table~\ref{tab:replication_results} shows the replication performance across all evaluated models.

\begin{table}[t]
\centering
\small
\resizebox{\columnwidth}{!}{
\begin{tabular}{l l c c c}
\hline
Setting & Model & Accuracy & Macro F1 & ROC-AUC \\
\hline
Imbalanced & RandomForest & 0.7688 & 0.6992 & 0.8923 \\
Imbalanced & XGBoost & 0.7896 & 0.7569 & 0.9162 \\
Imbalanced & AdaBoost & 0.7458 & 0.7135 & 0.8315 \\
\hline
SMOTE & RandomForest & 0.7792 & 0.7579 & 0.8924 \\
SMOTE & XGBoost & 0.8167 & 0.8020 & 0.9229 \\
SMOTE & AdaBoost & 0.6979 & 0.6962 & 0.8393 \\
\hline
\end{tabular}
}
\caption{Replication results across imbalance settings and models.}
\label{tab:replication_results}
\end{table}

Table 1 summarizes the performance of the replicated models under both imbalanced and SMOTE-balanced training settings. Across all runs, XGBoost achieved the strongest overall performance, having the highest accuracy, Macro F1, and ROC-AUC scores in both settings. For the imbalanced configuration, XGBoost achieved an accuracy of 0.7896 and a Macro F1 score of 0.7569. When SMOTE balancing was applied, performance improved further to 0.8167 accuracy and 0.8020 Macro F1. Random Forest also demonstrated a stable performance across both settings, while AdaBoost showed a lower overall accuracy and less table class-level behavior.

The confusion matrices show that home-win and away-win outcomes were consistently easier to classify than draw outcomes. In the imbalanced setting, Random Forest was able to correct identify 34 draw matches, while XGBoost improved draw prediction with 59 correctly classified draws. After applying SMOTE, draw prediction improved further across all models, with SMOTE-balanced XGBoost correctly classifying 82 draw matches.

The per-class metric further shows the difficulty of predicting draw outcomes. Home-win and away-win classes usually achieved stronger precision and recall scores, while draw-class F1 scores remained lower across most models. Although SMOTE improved draw recall in several cases, draw predictions still showed less stability than the other outcome classes.


\section{Gap Analysis vs Original Work}

Several differences were observed between the original paper’s reported results and our replication outputs. The original study reported XGBoost as the best model in the imbalanced setting and Random Forest as the strongest model under SMOTE balancing. In our replication, XGBoost achieved the best overall performance in both settings with substantially higher accuracy values. 

These gaps are likely caused by differences in datasets, preprocessing assumptions, feature engineering, and partially specified experimental details in the source paper. Our current implementation also uses synthetic/sample validation data rather than the exact dataset described in the original study. 

Despite the differences, the replication reproduced important trends from the paper. Ensemble and boosting methods consistently performed well, while draw outcomes remained the most difficult class to predict. SMOTE balancing improved draw recall in multiple settings but did not completely resolve the class imbalance challenges.


\section{Extension and Future Work}

As an extension to the original paper, we propose a draw-aware robustness enhancement focused on improving prediction performance for draw outcomes. Both the original study and the replication results showed that the draw class is more difficult to classify than home-win and away-win outcomes.

The proposed extension includes three components. First, uncertainty-based parity features would capture closely matched teams using differences in form or attack strength. Second, cost-sensitive training would assign higher penalties to draw misclassification. Third, probability calibration and a draw-sensitive decision threshold would improve confidence estimation for balance matches. 

Future experiments would compare the extended pipeline against the current Random Forest, XGBoost, and AdaBoost baselines using the same evaluation metrics. The main goal is to improve draw-class F1 by at least 0.05 while limiting overall accuracy loss to below 0.02. Future work may also explore temporal sequence models such as LSTMs or Transformers to capture match-history trends over time.


\section{Conclusion}

This project presented a reproducible machine learning pipeline for football outcome prediction using Random Forest, XGBoost, and AdaBoost models under both imbalanced and SMOTE settings. The replication successfully reproduced the overall workflow and showed that draw outcomes remain the most difficult class to predict. In addition, a draw-aware extension was proposed to improve robustness and future model performance on balanced matches.


\nocite{*}
\bibliography{CS5170_Project_Paper}

\end{document}